% !PRINTED 443
\weberpage{443}
But the right-hand side of this formula is entirely independent
of $m'$. The value of the left-hand side is therefore the same as if
$m' = 1$, and if one substitutes for the products the values from
(4), it follows that

\[
\left(\frac{mm'}{n}\right) : \left(\frac{m'}{n}\right) = \left(\frac{m}{n}\right) : \left(\frac{1}{n}\right),
\]

or

\[
8.\quad \left(\frac{mm'}{n}\right) = \left(\frac{m}{n}\right)\left(\frac{m'}{n}\right).
\]

Here it was at first assumed that $m$ and $m'$ are positive and
odd. But by 3. and 5., 8. still remains valid
when $m$ or $m'$ is negative or even, provided only that $m$ and $m'$
are relatively prime to $n$.

If we again take $m$ odd and positive, it now follows
from (4) and (9)

\weq{10}{\left(\frac{m}{n}\right) = (-1)^{\frac{(m-1)(n-1)}{4}} 2^{\frac{(m-1)(n-1)}{2}} \overset{\alpha}{\Pi}\overset{\beta}{\Pi} (\alpha - \beta).}

If we interchange $m$ with $n$ here, then in the
double product on the right-hand side every factor changes its sign;
everything else remains unchanged. But the number of these factors
is

\[
\frac{n-1}{2} \cdot \frac{m-1}{2}
\]

and from this it follows

\[
9.\quad \left(\frac{m}{n}\right) = (-1)^{\frac{(m-1)(n-1)}{4}} \left(\frac{n}{m}\right),
\]

\gesperrt{or} $\left(\dfrac{m}{n}\right)$ \gesperrt{is equal to} $\left(\dfrac{n}{m}\right)$\gesperrt{, if of the two
odd numbers} $m, n$ \gesperrt{at least one is of the form}
$4k+1$ \gesperrt{, and} $\left(\dfrac{m}{n}\right)$ \gesperrt{is opposite
 to} $\left(\dfrac{n}{m}\right)$\gesperrt{, if both numbers are of the
form} $4k+3$\gesperrt{.}

This famous theorem is known under the name of the \gesperrt{law of reciprocity}.
 The proof given here is due to
\gesperrt{Eisenstein}. With its help, and by 5. and
6., one can determine the value of the symbol $\left(\dfrac{m}{n}\right)$ very quickly, by
%CONT>>
